\chapter{Implementierung von Markov-Ketten in Python}
\label{sec:appendixMarkovPython}
\section{Python-Implementierung von Markov-Ketten}

Markov-Ketten lassen sich mit Python sehr elegant implementieren und visualisieren. Im Folgenden zeigen wir eine einfache Implementation, die es erlaubt, Markov-Ketten zu simulieren und deren Eigenschaften zu analysieren.

\begin{lstlisting}[language=Python, caption=Grundlegende Implementation einer Markov-Kette in Python]
import numpy as np
import matplotlib.pyplot as plt

# Einfache Markov-Ketten Implementierung
class MarkovChain:
    def __init__(self, transition_matrix, states=None):
        """
        Initialisiert eine Markov-Kette mit einer Übergangsmatrix.

        Args:
            transition_matrix: Eine quadratische Matrix als numpy-Array
            states: Optional - Liste mit Namen der Zustände
        """
        self.P = np.array(transition_matrix)
        n = self.P.shape[0]
        if states is None:
            self.states = [str(i) for i in range(n)]
        else:
            self.states = states

        # Überprüfen, ob die Übergangsmatrix gültig ist
        row_sums = self.P.sum(axis=1)
        if not np.allclose(row_sums, 1.0):
            raise ValueError("Die Zeilen der Übergangsmatrix müssen sich zu 1 summieren.")

    def simulate(self, steps, start_state=0):
        """
        Simuliert die Markov-Kette für eine bestimmte Anzahl von Schritten.

        Args:
            steps: Anzahl der Simulationsschritte
            start_state: Anfangszustand (als Index)

        Returns:
            Liste der besuchten Zustände (als Indizes)
        """
        n = self.P.shape[0]
        if start_state >= n:
            raise ValueError(f"Startzustand muss kleiner als {n} sein")

        # Simulation initialisieren
        states = [start_state]
        current = start_state

        # Für jeden Schritt
        for _ in range(steps):
            # Übergangswahrscheinlichkeiten für den aktuellen Zustand
            probs = self.P[current]

            # Nächsten Zustand zufällig auswählen
            current = np.random.choice(n, p=probs)
            states.append(current)

        return states

    def compute_stationary_distribution(self):
        """
        Berechnet die stationäre Verteilung der Markov-Kette.

        Returns:
            Stationäre Verteilung als numpy-Array
        """
        # Eigenwerte und Eigenvektoren der transponierten Matrix
        eigenvalues, eigenvectors = np.linalg.eig(self.P.T)

        # Finde den Eigenvektor zum Eigenwert 1
        one_eigenvector = eigenvectors[:, np.isclose(eigenvalues, 1.0)]

        # Normalisieren, damit die Summe 1 ergibt
        stationary = one_eigenvector[:, 0].real
        stationary = stationary / stationary.sum()

        return stationary
\end{lstlisting}

Mit dieser Klasse können wir nun einfach Markov-Ketten simulieren und analysieren. Zum Beispiel können wir unser Wetter-Beispiel implementieren:

\begin{lstlisting}[language=Python, caption=Anwendung der MarkovChain-Klasse auf das Wetter-Beispiel]
# Wetter-Markov-Kette
weather_matrix = [
    [0.8, 0.2],  # Sonnig -> Sonnig, Sonnig -> Regnerisch
    [0.4, 0.6]   # Regnerisch -> Sonnig, Regnerisch -> Regnerisch
]

weather_chain = MarkovChain(weather_matrix, states=["Sonnig", "Regnerisch"])

# Simulation für 20 Tage, beginnend mit sonnigem Wetter
simulation = weather_chain.simulate(20, start_state=0)
weather_states = [weather_chain.states[s] for s in simulation]

print("Wettersimulation für 20 Tage:")
print(weather_states)

# Stationäre Verteilung berechnen
stationary = weather_chain.compute_stationary_distribution()
print("\nStationäre Verteilung:")
for i, state in enumerate(weather_chain.states):
    print(f"{state}: {stationary[i]:.2f}")
\end{lstlisting}

Eine mögliche Ausgabe könnte sein:
\begin{verbatim}
Wettersimulation für 20 Tage:
['Sonnig', 'Sonnig', 'Sonnig', 'Regnerisch', 'Regnerisch', 'Sonnig', 
'Sonnig', 'Sonnig', 'Sonnig', 'Sonnig', 'Sonnig', 'Regnerisch', 
'Regnerisch', 'Regnerisch', 'Sonnig', 'Sonnig', 'Sonnig', 'Sonnig', 
'Sonnig', 'Sonnig', 'Sonnig']

Stationäre Verteilung:
Sonnig: 0.67
Regnerisch: 0.33
\end{verbatim}

Diese Python-Implementation erlaubt es uns, schnell verschiedene Markov-Ketten zu modellieren und zu simulieren. Besonders nützlich ist die Möglichkeit, die stationäre Verteilung zu berechnen und zu visualisieren, wie sich die Zustandsverteilung über die Zeit entwickelt.

\begin{myexercise}[Python-Simulation]
    Erweitern Sie den obigen Code, um die App-Nutzung aus unserem früheren Beispiel zu simulieren. Die Übergangsmatrix ist:
    \[
        P = \begin{pmatrix}
            0{,}7 & {,}3 \\
            {,}3  & {,}7 \\
        \end{pmatrix}
    \]

    Führen Sie eine Simulation für 100 Schritte durch und berechnen Sie die stationäre Verteilung. Vergleichen Sie die stationäre Verteilung mit dem analytischen Ergebnis aus dem früheren Beispiel.

    \begin{myanswer}
        Hier ist der Code für die Simulation der App-Nutzung:

        \begin{lstlisting}[language=Python]
# App-Nutzungs-Markov-Kette
app_matrix = [
    [0.7, 0.3],  # WhatsApp -> WhatsApp, WhatsApp -> Instagram
    [0.3, 0.7]   # Instagram -> WhatsApp, Instagram -> Instagram
]

app_chain = MarkovChain(app_matrix, states=["WhatsApp", "Instagram"])

# Simulation für 100 Nutzungen, beginnend mit WhatsApp
simulation = app_chain.simulate(100, start_state=0)

# Zählen, wie oft jede App genutzt wurde
whatsapp_count = simulation.count(0)
instagram_count = simulation.count(1)
print(f"WhatsApp wurde {whatsapp_count} mal genutzt")
print(f"Instagram wurde {instagram_count} mal genutzt")
print(f"Verhältnis WhatsApp/Instagram: {whatsapp_count/(whatsapp_count+instagram_count):.2f}")

# Stationäre Verteilung berechnen
stationary = app_chain.compute_stationary_distribution()
print("\nStationäre Verteilung:")
for i, state in enumerate(app_chain.states):
    print(f"{state}: {stationary[i]:.2f}")
\end{lstlisting}

        Die stationäre Verteilung ist $\pi = (0.5, 0.5)$, was genau dem analytischen Ergebnis entspricht. Bei einer langen Simulation sollte sich das Verhältnis der App-Nutzung diesem Wert annähern.
    \end{myanswer}
\end{myexercise}